% Generated draft: requires mathematical and source review.
\documentclass[11pt]{article}
\usepackage[margin=25mm]{geometry}
\usepackage{fontspec}
\setmainfont{Latin Modern Roman}
\usepackage{amsmath,amssymb,mathtools,unicode-math}
\setmathfont{Latin Modern Math}
\usepackage{enumitem,needspace,xurl,hyperref}
\setlength{\parindent}{0pt}
\setlength{\parskip}{5pt}
\setlength{\emergencystretch}{3em}
\newcommand{\N}{\mathbb{N}}
\newcommand{\Z}{\mathbb{Z}}
\newcommand{\Q}{\mathbb{Q}}
\newcommand{\R}{\mathbb{R}}
\newcommand{\C}{\mathbb{C}}
\newcommand{\F}{\mathbb{F}}
\newcommand{\A}{\mathbb{A}}
\newcommand{\PP}{\mathbb{P}}
\newcommand{\E}{\mathbb{E}}
\newcommand{\eps}{\varepsilon}
\newcommand{\dd}{\,\mathrm{d}}
\begin{document}
\section*{UnsolvedMath ID: 2306060}
\textbf{Source problem code:} catalogue entry
\subsection{Definitions and mathematical statement}
Let \(C\subset\C\) be a Jordan curve with bounded complementary domain \(\Omega_{\mathrm{in}}\) and unbounded complementary domain \(\Omega_{\mathrm{out}}\), considered together with \(\infty\). Suppose there are conformal maps
\[
f:\mathbb D\to\Omega_{\mathrm{in}},\quad
f(z)=z+\sum_{n\ge2}a_nz^n,
\]
\[
g_1:\mathbb D\to\Omega_{\mathrm{out}}\cup\{\infty\},\quad
g_1(z)=b_{-1}z^{-1}+b_0+\sum_{n\ge1}b_nz^n,\quad b_{-1}\ne0.
\]
The second map is meromorphic at \(0\), with a simple pole mapping \(0\) to \(\infty\); it is not an everywhere finite holomorphic disk map. The normalization of \(f\) in particular places \(0\) inside \(C\) and sets its conformal radius there to \(1\).

The source records that imposing \(b_{-1}=1\) forces \(C=\mathbb T=\{|z|=1\}\), by the area principle. Brannan's problem removes that normalization and asks what relationships can then be established between \(f\) and \(g_1\), including: asymptotics of their Taylor/Laurent coefficients; their boundary Lipschitz behavior on \(\mathbb T\); and related prime-end information when the boundary is allowed to be ``slightly non-Jordan.'' Boundary Lipschitz behavior may, for example, concern estimates of the form \(|f(\zeta)-f(\eta)|\le L|\zeta-\eta|\), and analogous estimates for \(g_1\), with the exact regularity target to be specified. Prime ends are the ideal boundary elements used in conformal compactifications of simply connected planar domains.

This is a request for relations or comparison theorems, not one specified coefficient identity or universal regularity claim. The source does not fix coefficient asymptotic scales, the desired boundary exponents or direction of implication, restrictions on \(b_{-1}\), or what non-Jordan boundaries and prime-end equivalences are permitted. Its Jordan setting and its proposed extension must therefore be kept distinct. The supplied OPEN-TRIAGE summary reports open status in the 2018 edition and no definitive resolution located.
\subsection{Short English statement}
Relate the normalized interior conformal map of a Jordan curve to its exterior map when the exterior leading Laurent coefficient is unrestricted. The source asks broadly about coefficients, boundary regularity, and non-Jordan prime-end extensions, without a single precise theorem target.
\subsection{Sources}
\begin{itemize}
\item Local UnsolvedMath registry, permanent ID \(2306060\), exact source code \texttt{catalogue entry}, ``Research Problems in Function Theory --- Problem 6.60.''
\item Hayman, \emph{Research Problems in Function Theory} (2018), Problem 6.60, attributed to D. A. Brannan; \url{https://arxiv.org/abs/1809.07200}. Source-TeX extraction; recorded access 2026-07-29.
\item The source invokes the area principle without a separate citation. The supplied OPEN-TRIAGE summary gives no authoritative current-resolution reference.
\end{itemize}
\end{document}
